\section{System Model}
\label{sec:model}

\subsection{Architecture}

We model a single-tier distributed service consisting of $N$
homogeneous worker nodes fronted by a load balancer. Requests arrive as
a Poisson process with aggregate rate $\Lambda = L \cdot N \cdot \mu$,
where $L \in (0, 1)$ is the \emph{load factor} and $\mu = 1/\tbase$ is
the per-node service rate. Each worker processes requests sequentially
(single-server) with deterministic service time $\tbase$ under normal
conditions, yielding an M/D/1 queueing model per node.

The load balancer implements weighted Power-of-Two-Choices (P2C): for
each request, it samples two candidates with probability proportional
to their weights $w_i \in [0, 1]$ and routes to the candidate with the
shorter queue depth.

\subsection{Fault Model}

A \emph{slow fault} on node $i$ multiplies its service time by a
slowdown factor $s_i > 1$:
\begin{equation}
  t_i = s_i \cdot \tbase
\end{equation}
We consider $f$ simultaneously faulty nodes out of $N$ total. Faults
are \emph{persistent}: once onset occurs, the slowdown remains until
explicit recovery (if any). We study five fault patterns:

\begin{itemize}
\item \textbf{Permanent}: constant $s$ after onset.
\item \textbf{Progressive}: $s(t) = 1 + (s_{\max}-1)(1-e^{-\beta t})$,
  gradually worsening.
\item \textbf{Fluctuating}: alternates between $s = s_{\text{peak}}$ and
  $s = 1$ with configurable on/off durations.
\item \textbf{Cascading}: faults appear sequentially on different nodes
  at staggered times.
\item \textbf{Recovering}: fault active for a fixed duration, then node
  returns to normal.
\end{itemize}

\subsection{Performance Metrics}

We evaluate using standard tail latency metrics:
\begin{itemize}
\item \textbf{P99/P999 latency}: 99th and 99.9th percentile of
  end-to-end request latency (including queueing delay).
\item \textbf{SLO violation rate}: fraction of requests exceeding the
  SLO target (50\,ms in our experiments).
\item \textbf{Goodput}: requests completing within SLO per second.
\item \textbf{Affected ratio}: fraction of requests served by faulty
  nodes (measures load balancer effectiveness).
\end{itemize}

\subsection{Minimax Regret}

To compare strategies across diverse scenarios, we use the \emph{minimax
regret} criterion. For strategy $j$ in scenario $i$, the regret is:
\begin{equation}
  R_{ij} = P99_{ij} - \min_k P99_{ik}
\end{equation}
The minimax regret of strategy $j$ is $\max_i R_{ij}$: its worst-case
gap from the best strategy in any scenario. A strategy with low minimax
regret is robust across operating conditions.
